\documentclass[10pt,a4paper]{amsart}
\usepackage[margin=19mm]{geometry}
\usepackage{amsmath,amssymb,mathtools}
\usepackage[scr=rsfs,cal=euler]{mathalfa}
\usepackage{xfrac}
\usepackage[unicode,hidelinks,bookmarks=false]{hyperref}
\newcommand{\ie}{i.e.\ }
\newcommand{\cf}{cf.\ }
\newcommand{\resp}{respectively\ }
\newcommand{\Subsection}{\subsection}
\providecommand{\nobreakdash}{\nobreak\hbox{-}}
\setlength{\emergencystretch}{2em}
\multlinegap0pt
\usepackage{bm}
\usepackage{bbm}
\usepackage{dsfont}
\usepackage{xspace}
\usepackage{cancel}
\usepackage{mathtools}
\usepackage{amsthm}
\makeatletter
\@ifundefined{thm}{\newtheorem{thm}{Thm}}{}
\@ifundefined{lem}{\newtheorem{lem}[thm]{Lemma}}{}
\makeatother
\makeatletter
\expandafter\ifx\csname claimproof\endcsname\relax
\newenvironment{claimproof}[1][\claimproofname]{\begin{proof}[#1]\renewcommand*{\qedsymbol}{\(\square\)}}{\end{proof}}
\fi
\makeatother
\title[A minimum-degree threshold for colour-biased Hamilton cycles]{A minimum-degree threshold for colour-biased Hamilton cycles in hypergraphs}
\author{Natalie Behague, Felix Christian Clemen, Joseph Hyde, Natasha Morrison}
\date{Source: arXiv:2607.29628v2 (CC BY 4.0)}
\begin{document}
\maketitle

\noindent\emph{Source and licence.} A minimum-degree threshold for colour-biased Hamilton cycles in hypergraphs. Version arXiv:2607.29628v2 (CC BY 4.0), distributed under the Creative Commons Attribution 4.0 International licence, as recorded in the arXiv licence metadata for this exact version and shown on the abstract page https://arxiv.org/abs/2607.29628v2. Full rights notice: licenses/arxiv-2607.29628-CC-BY-4.0.txt.\par\smallskip

\noindent\textit{Editorial scope.} Counted unit: the Lem whose source label is \texttt{annoying} in the article (source0.tex lines 595--599, source label \texttt{annoying}), delivered together with the article's own complete original proof of it (source0.tex lines 601--629, one \texttt{proof} environment of 29 lines). No mathematics is written, repaired or re-derived: statement and proof are the article's own text line for line. Required context retained uncounted: none. Every definition, display and label of those lines is retained unchanged. Omitted from this package and not redistributed: 1--594, 600--600, 630--1285. Nothing inside a retained excerpt is omitted or altered: every retained body is delivered exactly as the article's lines are written, so no omission receipt of any kind accompanies this package. Citations: the retained excerpts carry no citation command at all, so no bibliography entry of the article is transcribed and no reference is redistributed. Reference adaptation: none was needed and none was made; every reference inside the delivered unit resolves inside it, so no unresolved reference or citation is produced. Printed slips kept literal: no printed word, symbol, hypothesis or number of the counted statement or of its proof was altered. Layout and class: the article's own class is \texttt{article} with packages including fullpage, amssymb, amsmath, amsthm, amsfonts, comment, mathtools, thm-restate, tikz, subcaption, hyperref, cleveref, xcolor, enumerate; the wrapper is the lane's pinned \texttt{amsart} editorial wrapper at 10pt on A4, which restates the theorem-like environments the retained text needs and adds the article's own macro definitions that the retained text needs (none), with extra wrapper packages (bm, bbm, dsfont, xspace, cancel, mathtools). The delivered numbering is the unit's own. Assets: no retained span contains a figure environment, a graphics-inclusion command, a PDF-page inclusion, a table or a TikZ picture; no image file is copied into this package, the package declares no asset dependency and no image file is redistributed with it. Licence: version arXiv:2607.29628v2 of this article is distributed under the Creative Commons Attribution 4.0 International licence, as recorded in the arXiv licence metadata for this exact version and shown on the abstract page https://arxiv.org/abs/2607.29628v2. A full rights notice is delivered at licenses/arxiv-2607.29628-CC-BY-4.0.txt. Characters: every retained excerpt is pure ASCII, so the delivered engine, XeLaTeX with its default Latin Modern text font, reports no missing character.
\begin{lem}
\label{annoying}
A non-agreeable component on $n$
vertices has at most $6n$ edges. 
\end{lem}

\begin{proof}
Let $G$ be a non-agreeable component, and fix a labelling
$w:E(G)\to\{-2,0,2\}$ witnessing this. Choose a longest path
$P=v_1v_2\cdots v_k$, and let $a_i := w(v_iv_{i+1})$. If $k \le 3$, then $G$ is a star or triangle and $e(G) \le n$. Therefore we will assume $k \ge 4$. 

If $P$ contains two consecutive edges with weight 0, then (as every copy of $P_3$ has weight 0) every edge of $P$ has weight 0. Now any edge incident to $P$ also has weight 0 (as it forms a 3 edge path with two edges already labelled 0). By induction on the distance of an edge from $P$, the same argument shows
that every edge in $G$ has weight $0$, a contradiction. So $P$ contains an edge with non-zero weight. So, up to permuting the values $2,0,-2$,
the weights along $P$ are $2,0,-2,2,0,-2,\ldots$. 

Let $3\le i\le k-2$. We show that $d(v_i) \le 4$. First suppose that
$x \in X := V(G) \setminus P$ is adjacent to $v_i$. Then both
$v_{i-2}v_{i-1}v_ix$ and $xv_iv_{i+1}v_{i+2}$ are paths of length three.
Hence $w(v_ix)=-(a_{i-2}+a_{i-1})=a_i$ and $w(v_ix)=-(a_i+a_{i+1})=a_{i+2}$, a contradiction as $a_i\ne a_{i+2}$. Now let
$e=v_iv_j$ be a chord with $|i-j|\ge 3$. Here the paths
$v_{i-2}v_{i-1}v_iv_j$ and $v_jv_iv_{i+1}v_{i+2}$ give a contradiction as above. So $N(v_i) \subseteq \{v_{i-2},v_{i-1},v_{i+1},v_{i+2}\}$, as required. By maximality of $P$, we obtain in addition that each vertex of $X$ is necessarily a neighbour of $v_2$ or $v_{k-1}$. 

Note that $v_1$ and $v_k$ cannot have a neighbour outside $P$, as otherwise $P$ is not the longest path. Thus by the observations above, $N(v_1) \subseteq \{v_2,v_3,v_{k-1},v_k\}$ and so $d(v_1) \le 4$ and correspondingly $d(v_k) \le 4$. 

Now consider $W = N(v_2) \setminus V(P)$. We will
show that vertices in $W$ are only adjacent to $v_2$. Suppose not and let $x \in W$ and $y \ne v_2$ be a neighbour of $x$. Note that $y \ne v_1$ as $x \not\in N(v_1) \subseteq V(P)$. Moreover $y \ne v_3$: if $3\le k-2$ then $x \not\in N(v_3) \subseteq V(P)$, and if $3=k-1$ then $y = v_3$ would result in the path $v_1v_2xv_3v_4$, which is longer than $P$. Thus both $v_1v_2xy$ and $v_3v_2xy$ are copies of $P_3$ and have weight 0. In particular, $w(v_1v_2) = w(v_2v_3)$, which is a contradiction as $a_1 \ne a_2$.
By a corresponding argument, vertices in $W' = N(v_{k-1})\setminus V(P)$ are only adjacent to $v_{k-1}$.




We conclude that every edge of $G$ contains a vertex of $P$. Hence, we have a crude bound of
$$e(G) \le \sum_{i=1}^k d(v_i) \le 4(k-2) + 2n  \le 6n.$$

\end{proof}

\end{document}
